% Figure: rectangular (Cartesian) gain pattern of a uniformly illuminated
% line/aperture array, gain in dB versus angle. The array factor magnitude
% for N elements is |sin(N psi/2) / (N sin(psi/2))| with psi = (2 pi d/lambda) sin theta.
% Here N = 8, d = lambda/2, so psi = pi sin theta. Plotted in dB to show the
% main lobe, the first nulls, and the -13 dB first sidelobes.
\begin{figure}[htbp]
\centering
\begin{tikzpicture}
\begin{axis}[
  width=12cm, height=6.5cm,
  xlabel={angle from broadside $\theta$ (degrees)},
  ylabel={normalised gain (dB)},
  xmin=-90, xmax=90, ymin=-40, ymax=2,
  domain=-90:90, samples=600,
  axis lines=left,
  xtick={-90,-60,-30,0,30,60,90},
  ytick={0,-10,-13,-20,-30,-40},
  enlargelimits=false,
  restrict y to domain=-40:5,
]
  % AF magnitude, N=8, psi = pi sin theta (theta in degrees -> sin)
  % 20 log10 |sin(N psi/2)/(N sin(psi/2))|, with a small floor to avoid log(0)
  \addplot[engblue, very thick]
    { 20*log10( abs( sin(4*180*sin(x)) ) / ( 8*abs(sin(0.5*180*sin(x))) + 0.0001 ) + 0.0001 ) };
  % first-sidelobe reference line at -13 dB
  \addplot[engred, dashed, thick, domain=-90:90] {-13};
  \node[engred, font=\scriptsize, anchor=west] at (axis cs:55,-11) {$-13$ dB sidelobe};
\end{axis}
\end{tikzpicture}
\caption[Gain pattern of an eight-element broadside array]{The
normalised gain pattern of an eight-element uniform broadside array with
half-wavelength spacing, plotted in decibels against angle from
broadside. The narrow main lobe at $\theta = 0$ is flanked by nulls and
by sidelobes; the first sidelobe of a uniformly fed array sits about
$13\,\si{\decibel}$ below the main-lobe peak, the price paid for uniform
illumination. Tapering the element amplitudes lowers the sidelobes at
the cost of a wider main lobe.}
\label{fig:vol04ch10:gain-pattern}
\end{figure}
